\subsection{The Cost-Modification Mechanism}
\label{sec:penalization}

The framework fixes neither which environmental variables alter the network nor the magnitude of their effect. It exposes a \emph{mechanism} that receives environmental layers and returns a new cost for each edge; selecting the layers and setting their parameters belongs to whoever applies the method. This section describes that mechanism.

\subsubsection{Scenarios}

The network is evaluated in two states. In the \textbf{typical scenario}, each edge weight ($W_{base}$) is the travel cost under normal conditions — the segment length divided by the mean speed adopted for the mode under analysis. In the \textbf{conditioned scenario}, weights are recomputed from the environmental layers ($C$), producing $W_{cond}$ for each edge.

Comparing the two states isolates the effect of the disruption: any difference in accessibility indicators is attributed to the cost change introduced by the environmental layers.

\subsubsection{Cost modifiers}

The core element of the mechanism is the \emph{cost modifier}. Each modifier is a function that takes the value of a layer at a point on the edge and returns a multiplicative factor:

\begin{equation}
  f(v) \rightarrow \mathbb{R}^{+}
\end{equation}

where $v$ is the value sampled from the layer at that point. The factor $f$ may exceed 1 (penalization), fall below 1 (attenuation), or be infinite (interdiction). The same mechanism covers all three cases, which is why the framework refers to \emph{cost modifiers} rather than penalties alone.

A modifier never receives the whole edge: it is applied to values sampled along it. Sampling strategy and factor aggregation are declared separately, so phenomena with different spatial distributions are handled by the same mechanism.

\subsubsection{Sampling}

Two sampling strategies are available, and the choice trades cost for fidelity. \textbf{Midpoint} queries the layer at the edge midpoint and applies a single factor to the entire edge, which is fast but inaccurate when the edge crosses a boundary. \textbf{Segments} places $n$ equidistant points along the edge, including both endpoints, and reads the layer at each; the value of $n$ is a parameter, defaulting to four.

The distinction matters most for long edges, which are common in peripheries where blocks are large. A single edge can span several hundred meters and cross more than one thermal or topographic zone. Sampling only its midpoint may place the entire segment in the wrong class.

\subsubsection{Reading the layers}

How a layer is queried depends on the nature of the data.

\textbf{Continuous data} (rasters) are sampled at the chosen points. Each point takes the value of the pixel containing it. Points landing on a nodata pixel leave the cost unchanged at that location rather than propagating a missing value into the cost, so an incomplete raster degrades coverage instead of corrupting the network.

\textbf{Discrete data} (vector polygons) are queried by intersection. The edge is partitioned at every point where a polygon boundary crosses it, and each resulting sub-segment receives the factor of the polygon covering it. Sub-segments outside any polygon keep the original cost. Where polygons overlap, the \emph{largest} factor among them applies to the overlap, so a region mapped by two hazard layers is governed by the more restrictive one.

\subsubsection{Aggregation and interdiction}

The factors obtained along an edge must be reduced to a single value for the whole edge. The framework offers two reductions:

\begin{itemize}
  \item \textbf{Maximum}: adopts the most restrictive factor observed along the segment, treating the edge by its worst case. This is conservative and appropriate when the sample is small.
  \item \textbf{Length-weighted mean}: each sub-segment pays its own factor, and the mean is weighted by extent. This is more faithful on long edges crossing heterogeneous regions.
\end{itemize}

Infinite factors receive separate treatment. Assigning $f = \infty$ to a sub-segment declares it impassable: the cost of any path crossing it becomes infinite, and routing prefers alternatives. Interdiction is not a high cost but a statement that crossing is not viable in that scenario — the appropriate treatment for phenomena that block the road rather than delay it, such as a landslide burying the segment. When no alternative exists, the cell reaches no destination at all, a situation recorded as access loss.

\subsubsection{Composition}

When more than one environmental layer affects the same edge, modifiers are applied in sequence and the result accumulates multiplicatively:

\begin{equation}
  W_{cond} = W_{base} \times f_1 \times f_2 \times \cdots \times f_n
\end{equation}

Each step receives the cost produced by the previous one. Because multiplication commutes, the order does not change the numeric result, but it changes how the process reads, which matters when attributing an accessibility change to a specific layer.

This composition admits two arrangements. In \emph{simple rules}, each modifier queries one layer independently and chaining happens across rules. In \emph{composite rules}, several layers are queried at the same point on the edge and the factor is computed jointly, allowing one layer's effect to depend on another's value.

The distinction is not cosmetic. In the simple arrangement, an attenuation by vegetation would apply to any tree-lined segment, including where there is no heat at all — dense canopy next to a shaded, cool street would still receive a discount. In the composite arrangement, vegetation attenuates only where the thermal layer exceeds a threshold at that same point, so the two factors are evaluated jointly rather than independently averaged. A composite rule thus samples all its layers at identical locations before aggregating, preserving the spatial correspondence between their values.

\subsubsection{Recalculation and indicators}

With conditioned costs defined, a least-cost path algorithm runs over the graph, producing travel-time matrices for both scenarios. From these, the indicators described in Section~\ref{subsec:territorial_modelling} are computed: the mean time to reach the $k$ nearest POIs of each category (PTh), the inequality of its distribution (index G), and the share of the population served within a threshold (F15).

The scenario comparison adds per-cell time variation and \emph{access loss}: cells reaching no POI under the conditioned scenario receive infinite time and are accounted for separately. They are not an extreme case of high time but the absence of access, and treating them as a numeric value would distort any mean or inequality index computed over the remainder.
